% Erzeugen: latexmk -pdf main.tex   (empfohlen)
% Alternativ: pdflatex main.tex && pdflatex main.tex

\documentclass[11pt,a4paper]{report}

\usepackage[T1]{fontenc}
\usepackage[utf8]{inputenc}
\usepackage{lmodern}
\usepackage[ngerman]{babel}
\usepackage{microtype}
\usepackage{amsmath}

\usepackage[a4paper,margin=2.5cm]{geometry}

\usepackage{graphicx}
\usepackage{booktabs}

\usepackage[hidelinks]{hyperref}

% Doppelte PDF-Anker auf Titelseite und Abstract vermeiden (report-Klasse).
\hypersetup{pageanchor=false}

\title{Dekzentrale Predictionsmarkt-DApp\\\large Technische Dokumentation (\texttt{blch-voting-dapp})}
\author{Ihr Name(n)\\\small Institution oder Lehrveranstaltung}
\date{\today}

\begin{document}

\maketitle

\begin{abstract}
Dieses Dokument beschreibt die im Repository \texttt{blch-voting-dapp}
umgesetzte Demonstrations-DApp: einen vereinfachten, poolbasierten Predictionsmarkt
auf Ethereum (Solidity), angebunden an ein Nuxt-Browser-Frontend mit
MetaMask. Nutzerinnen und Nutzer beziehen Demonstrations-Token \emph{VOTE},
eröffnen Märkte mit zwei Ausgängen (im Frontend fest \emph{YES}/\emph{NO}),
staken auf einen Ausgang und erhalten nach Auflösung durch eine
Admin-Adresse anteilige Gewinne aus dem gesamten Markt-Pool. Kernbegriffe sind
ERC-20, Wallet-gestützte Transaktionen, On-Chain-State und ein
\emph{pull}-basiertes Belohnungsmodell ohne zentrale Geschäftslogik auf einem Server.
\end{abstract}
\hypersetup{pageanchor=true}

\tableofcontents

\clearpage

\section{Einleitung}
Diese Arbeit beschreibt eine reduzierte Predictionsmarkt-DApp. Die Anwendung
verbindet einen ERC-20-Demonstrationstoken, marktbasierte Einsaetze auf diskrete
Ausgaenge und eine administrative Aufloesung nach Fristende. Im Mittelpunkt steht
damit kein allgemeines Blockchain-Geruest, sondern ein klar abgegrenzter
Kernfluss mit Tokenverteilung, Markterstellung, Einsatz und Auszahlung.

Als Beispiel dient ein Predictionsmarkt. Das Beispiel stellt fachliche und
technische Fragen kompakt dar. Die Teilnahme erfordert keine tiefe Vorarbeit.
Teilnehmende koennen direkt an einem binären Markt mitwirken. Gleichzeitig sind
Ergebnisfindung, Vertrauensannahmen und Informationsasymmetrien zentrale
Aspekte.

Das übergeordnete Ziel besteht darin, eine durchgängig lauffaehige DApp
bereitzustellen, deren Aufbau, Vertragslogik und Kopplung zur Oberfläche
nachvollziehbar bleiben. Konkret adressiert die Implementierung folgende Leitfragen:
\begin{itemize}
  \item Wie lassen sich Markterstellung, Teilnahme und Abrechnung so in Smart
    Contracts formulieren, dass der Zustand vollständig on-chain nachvollziehbar
    bleibt?
  \item Wie gestaltet sich die Kopplung zwischen Browser-Oberflaeche, Wallet und
    JSON-RPC-Endpunkt, sodass typische Nutzeraktionen (Genehmigen, Staken,
    Auflösen, Beanspruchen) robust abgebildet werden?
  \item Welche Sicherheits- und Vertrauensannahmen sind für eine Demo akzeptabel,
    und wo müsste ein produktives System nachrüsten?
\end{itemize}

Ziel dieses Berichts ist eine kurze, fachlich geschlossene technische
Dokumentation im Stil einer Projektarbeit. Sie beschreibt die Lösung, analysiert
die wichtigsten Designentscheidungen und leitet daraus die Grenzen und das
Verbesserungspotenzial der Implementierung ab.

Kapitel~\ref{chap:theorie} beschreibt die Methodik und die technische Grundidee.
Kapitel~\ref{chap:kern} analysiert den Kernablauf der DApp und den
Funktionsumfang. Kapitel~\ref{chap:sicherheit} diskutiert Grenzen und
Verbesserungspotenzial. Zusammen ergibt sich daraus ein kurzer, geschlossen
argumentierter Projektbericht.

Die Arbeit ist als Demonstrationssystem konzipiert und nicht für den Einsatz mit
realen Daten vorgesehen. Die folgenden Punkte sind daher keine nachträglichen
Defizite, sondern bewusst gesetzte Grenzen des Projektumfangs.

Als Sicherheiten-Token dient ausschliesslich der eigene ERC-20-Token \emph{VOTE} mit
einmaligem Faucet. Jede Adresse kann sich einmalig einen Betrag zuweisen. Damit
werden weder natives Ether noch ein Stablecoin als Einsatzmittel verwendet.
Preisrisiko, Compliance und Brueckenlogik bleiben damit ausserhalb des Systems.
Brueckenlogik bezeichnet hier die Komplexitaet, die sich bei Abhaengigkeiten zu
anderen Systemen oder zum Mainnet ergibt.

Die Aufloesung erfolgt durch eine beim Deployment gesetzte \emph{Admin}-Adresse.
Diese Rolle bestimmt nach Fristende den gewinnenden Ausgang. Das System enthält
keine Chainlink-Feeds, keine Reporter-Jury und keine Einspruchsinstanz. Chainlink-
Feeds sind externe Datenquellen, die in einem realen Anwendungsfall zur Aufloesung
eines Markts noetig sein koennen.

der Vertrag \texttt{PredictionMarket} erlaubt formal mehrere Ausgaenge, sofern
die Validierung in \texttt{createMarket} erfuellt ist. Die Nuxt-Oberfläche
vereinfacht die Nutzerfuehrung, indem sie neue Maerkte derzeit auf die beiden
festen Ausgaenge \emph{YES} und \emph{NO} festlegt. Die Benutzeroberfläche bildet
damit bewusst einen binären Kernfall ab.

Im weiteren Verlauf bezeichnet der Begriff \emph{Predictionsmarkt} die konkret
implementierte Variante: Ein Markt besitzt eine textuelle Frage, einen
Aufloesungskontext zur operationalen Klaerung, eine absolute Deadline, eine
Menge benannter Ausgaenge und pro Ausgang einen Token-Pool. Solange die Deadline
nicht überschritten ist, koennen Adressen Token in die Pools einbringen
(\emph{staken}). Nach Fristende wählt die Admin-Adresse den gewinnenden Ausgang.
Gewinner beanspruchen mit einer eigenstaendigen Transaktion einen anteiligen
Anteil am gesamten Markt-Pool gemaess der im Vertrag implementierten Formel.

Die aus den Poolverhaeltnissen abgeleiteten Prozentwerte sind in diesem Kontext
\textbf{einsatzgewichtete Anteile}. Sie dienen dem Nutzenden als Feedback fuer
einen freiwillig platzierten Einsatz. Es handelt sich somit nicht um eine exakte
Vorhersage eines realen Ereignisses.

\section{Hauptkapitel}
\label{chap:kern}

Das System implementiert den Kernablauf eines Predictionsmarkts. Eine Adresse legt
einen Markt mit Frage, Auflösungskontext, Laufzeit und Ausgängen an. Andere
Adressen beziehen Token über das Faucet, genehmigen den Marktvertrag und staken
auf einen Ausgang. Solange die Frist läuft, sammeln sich Einsätze in den Pools.
Nach Fristende setzt die Admin-Adresse den gewinnenden Ausgang. Gewinner fordern
ihren Anteil in einer eigenen Transaktion an.

Die Benutzeroberfläche reduziert neue Märkte auf \emph{YES} und \emph{NO} für
die Demonstration. Der Vertrag unterstützt formal mehrere Ausgänge. Die
Implementierung bleibt dabei transparent und nachvollziehbar.

Die Implementierung beruht auf drei Entscheidungen. Das Faucet-Modell senkt die
Einstiegshürde, da keine reale Wertposition nötig ist. Der ERC-20-Ablauf mit
\texttt{approve} und \texttt{transferFrom} führt zu einem zweistufigen Zahlungs-
fluss. Die Admin-basierte Auflösung schafft einen zentralen Vertrauenspunkt.

Gas- und Benutzungsaspekte sind reduziert. Die Anwendung zeigt nur die für den
Kernpfad notwendigen Transaktionen und vermeidet zusätzliche on-chain Schleifen
oder Sammelauszahlungen. Das Ergebnis ist ein schmaler, leicht verständlicher
Funktionsumfang.

\chapter{Systemarchitektur und Technologie}
\label{chap:arch}

\section{DApp-Kernkonzepte und ihre Umsetzung im Projekt}

Die Lehrveranstaltung verlangt typischerweise die sichere Anwendung
grundlegender DApp-Bausteine: eine öffentliche Blockchain als Vertrauensanker,
in Solidity geschriebene Smart Contracts, ein Browser-Frontend und eine Wallet,
die Transaktionen signiert. Im Merkblatt werden beispielhaft Web3.js und MetaMask
genannt; funktional äquivalent setzt dieses Projekt auf \textbf{ethers.js in der
Version 6} und ebenfalls MetaMask als injizierten Ethereum-Provider
(\texttt{window.ethereum}).

Das Frontend erzeugt für jeden benötigten Vertrag eine \texttt{Contract}-Instanz
mit ABI und Deployment-Adresse. Lesezugriffe nutzen den Provider; schreibende
Aufrufe laufen über einen \texttt{Signer}, der von der verbundenen Wallet
kommt. Damit ist die kryptographische Autorisierung klar von der bloßen
Datenabfrage getrennt: Nur wenn Nutzerinnen eine Transaktion bestätigen, ändert
sich der On-Chain-Zustand.

\section{Schichtenmodell und Datenflüsse}

\begin{sloppypar}
Die Architektur lässt sich als dreischichtiges Modell beschreiben, das von der
Präsentation bis zur persistenten Chain reicht. In komprimierter Lesart ergibt
sich folgender Datenfluss: Die Nuxt-Oberfläche ruft über \texttt{ethers.js}~v6
die MetaMask-Provider-Schnittstelle auf; MetaMask signiert Transaktionen und
sendet sie an einen JSON-RPC-Endpunkt (lokal Hardhat, öffentlich z.\,B.\
Sepolia). Lesende Aufrufe treffen dieselben Verträge \texttt{VoteToken} und
\texttt{PredictionMarket} ohne Wallet-Bestätigung, sofern die Anwendung einen
reinen Provider-Pfad nutzt---im vorliegenden Frontend dominiert jedoch der
Signer-Pfad für Einheitlichkeit.
\end{sloppypar}

\paragraph{Präsentationsschicht.}
Das Verzeichnis \path{frontend/} enthält die Nuxt-Anwendung. Die zentrale
Oberfläche liegt in \path{frontend/app/app.vue}; wiederkehrende Logik ist in
Composables ausgelagert, insbesondere \path{useWallet.ts} für Verbindung, Kette
und Signer sowie \path{useContracts.ts} für Vertragsinstanzen, Aggregationslese-
routinen und Transaktionshüllen.

\paragraph{On-Chain-Logik.}
Die Verträge \path{contracts/VoteToken.sol} und
\path{contracts/PredictionMarket.sol} bilden die fachliche Wahrheit. Sie werden
mit Hardhat kompiliert, getestet und über \path{scripts/deploy.ts} ausgerollt.
Die Trennung zwischen Token und Markt folgt dem Prinzip der getrennten
Zuständigkeiten: Der Token kapselt Ausgabe und Übertragung nach ERC-20; der
Marktvertrag kapselt Geschäftsregeln und hält aggregierte Pools.

\paragraph{Build- und Integrationspfad.}
Nach dem Deployment schreibt \path{scripts/deploy.ts} die Adressen nach
\path{deployments/<netzwerk>.json}. Das Skript \path{scripts/export-frontend.ts}
überträgt ABIs und Adressen in \path{frontend/app/contracts/}, sodass die
Oberfläche ohne manuelles Kopieren mit der jeweiligen Deployment-Generation
übereinstimmt. Dieser Schritt ist leicht zu vergessen, aber für reproduzierbare
Demos unverzichtbar.

\section{Abwesenheit eines zentralen Anwendungsservers für die Geschäftslogik}

Marktlisten, Poolstände und Auszahlungsberechtigungen werden nicht in einer
Datenbank verwaltet, sondern aus der Chain gelesen. Ein klassischer Backend-Dienst
könnte höchstens als Cache, Indexer oder für Metadaten dienen; im vorliegenden
Repository ist bewusst darauf verzichtet, um die DApp-Eigenschaft nicht zu
relativieren. Die Ausnahme bildet die semantische Rolle der Admin-Adresse: Sie ist
kein Server, aber ein zentrales Vertrauenselement auf Protokollebene, weil sie
die Auflösungsentscheidung trifft (vgl. Kapitel~\ref{chap:sicherheit}).

\section{Netzwerke, Konfiguration und Betriebsmodi}

\paragraph{Lokale Entwicklung.}
Für den Alltag im Kurs eignet sich ein Hardhat-Node mit Chain-ID \texttt{31337}
und RPC-Endpunkt \texttt{http://127.0.0.1:8545}. MetaMask muss dieses Netzwerk
als benutzerdefinierte Chain eingetragen haben, damit Transaktionen nicht an das
falsche Netzwerk gesendet werden.

\paragraph{Öffentliches Testnet.}
Das Projekt unterstützt Sepolia-Deployments über Umgebungsvariablen wie
\texttt{SEPOLIA\_RPC\_URL} und \texttt{SEPOLIA\_PRIVATE\_KEY}. Private Schlüssel
dürfen niemals im Versionskontrolsystem, in CI-Logs oder in Container-Layern
landen; stattdessen sind Geheimnisse über die jeweilige Laufzeitumgebung zu
injizieren.

\paragraph{Konsistenz zwischen Frontend und Chain.}
Die Adressen in \path{frontend/app/contracts/addresses.json} müssen exakt zu der
Chain passen, gegen die MetaMask verbunden ist. Ein häufiger Fehler im Praktikum
ist die Kombination aus lokalem Node, aber noch ausstehendem
\texttt{export:frontend} nach einem frischen Deploy, oder ein Sepolia-Frontend mit
Hardhat-Adressen.

\chapter{Smart Contract}

\section{Datenmodell}
Wichtige Speichervariablen, Aufzählungen und Zustandsübergänge.

\section{Kernfunktionen}
Lebenszyklus: Deployment, Einrichtung von Vorschlag oder Wahl, Abstimmung,
Auszählung sowie ggf.\ Pausen- oder Notfallpfade.

\section{Gas und Erweiterbarkeit}
Kostenüberlegungen, verwendete Muster (oder explizit \emph{kein}
Upgrade-Pfad) sowie Hinweise zum Deployment.

\chapter{Frontend, Nutzerrollen und Deployment}

\section{Frontend (Nuxt)}
Die Haupt-UI liegt in \path{frontend/app/app.vue}:
\begin{itemize}
  \item Wallet: Verbinden, Anzeige von Adresse, Chain-ID,
    \emph{VOTE}-Saldo und \emph{Allowance} gegenüber \texttt{PredictionMarket}.
  \item Nicht-Admin: Faucet-Button (\texttt{claimFaucet}), Formular zur
    Markterstellung mit Dauer und Resolution Context---Outcomes sind fest YES/NO.
  \item Marktkarten: Status-Badge, Pools, angezeigte Prozentwerte, Staken mit
    vorherigem Approve für den jeweiligen Betrag.
  \item Admin-Adresse (aus exportierter Konfiguration): Auflösungsbuttons pro
    Markt nach Fristablauf (\texttt{resolveMarket}).
\end{itemize}
Die Composable-Datei \path{useContracts.ts} kapselt alle Lesevorgänge
(\texttt{getMarkets}-Schleife über Märkte und Outcomes) und Transaktionsaufrufe.

\section{Ablauf für Endnutzerinnen}
\begin{enumerate}
  \item Wallet verbinden, Netzwerk lokal oder Sepolia wählen (MetaMask-Einträge
    wie in der README beschrieben).
  \item \emph{VOTE} einmalig fauceten und ggf.\ Saldo aktualisieren.
  \item Markt eröffnen oder bestehenden wählen, \emph{Approve} für gewünschten
    Stake-Betrag, danach YES oder NO staken (solange Zeitfenster und nicht
    aufgelöst).
  \item Nach Auflösung und eigenem Gewinn \texttt{claimReward} auslösen.
\end{enumerate}

\section{Build- und Betriebsszenarien}
\subsection{Lokal (Entwicklung)}
Typische Reihenfolge aus der README:
\texttt{npm run chain} (Terminal~1); \texttt{npm run deploy:local} sowie
\texttt{npm run export:frontend} (Terminal~2); \texttt{cd frontend \&\& npm run dev}
(Terminal~3).

\subsection{Docker \& CI}
Ein Docker-Image bündelt das gebaute Frontend und die exportierten Contract-JSONs,
nicht jedoch Private Keys oder Deployment-Skripte. GitHub Actions
(\path{.github/workflows/}) können das Image erzeugen und z.\,B.\ nach GHCR pushen---
Details im Repo \path{Dockerfile} und \path{docker-entrypoint.sh}.

\section{Schnittstelle zur Chain}
Deployed Adressen und ABIs liegen unter \path{frontend/app/contracts/} (u.\,a.\
\texttt{addresses.json}, Artefakte). Frontend und Dokumentation müssen zur
gleichen Deployment-Generation passen (\texttt{export:frontend} nach jedem Deploy).

\section{Diskussion}
\label{chap:sicherheit}

Predictionsmärkte sind nutzerfreundlich, weil sie keine tiefen technischen
Vorkenntnisse erfordern. Damit bieten sie eine niedrige Einstiegsschwelle.
Gleichzeitig zeigen sie, dass Ergebnisfindung, Transparenz und
Informationsasymmetrien zentrale Fragen sind. Die DApp bildet diesen Spannungs-
bogen in vereinfachter Form ab.

Die Implementierung ist ein Demonstrationssystem und nicht produktionsreif.
Die Auflösung erfolgt admin-basiert. Eine zentrale Instanz bestimmt das Ergebnis
nach Fristende und erzeugt so einen Vertrauenspunkt. Im vorliegenden System
fehlen ein dezentrales Oracle- oder Streitbeilegungsmodell, sekundärer Handel
und weitergehende Absicherungen gegen missbräuchliche Markterstellung.

Auch Gas- und Benutzungsaspekte bleiben reduziert. Die Anwendung zeigt nur die
für den Kernfluss notwendigen Transaktionen und vermeidet zusätzliche
Komplexität. Das begrenzt Funktionsumfang und Realitätsnähe.

Ein dezentrales Oracle- oder Streitbeilegungsmodell reduziert die Abhängigkeit
von einer einzelnen Admin-Adresse. Nutzer können Ergebnisse vorschlagen. Die
Plattform vermittelt bei Konflikten oder Anfechtungen. Solche Konzepte finden
sich in produktiven Protokollen wie Polymarket.

Weitere Erweiterungen betreffen Marktvielfalt, Auswertungen zur Gas-Nutzung und
eine klarere Trennung zwischen UI-Vereinfachungen und on-chain-Allgemeingültigkeit.
Die aktuelle Lösung bildet einen didaktisch klaren Kern ab, der technisch
nachvollziehbar bleibt.

Ein dezentrales Oracle- oder Streitbeilegungsmodell reduziert die Abhängigkeit
von einer einzelnen Admin-Adresse. Nutzer können Ergebnisse vorschlagen. Die
Plattform vermittelt bei Konflikten oder Anfechtungen. Solche Konzepte finden
sich in produktiven Protokollen wie Polymarket.

Weitere Erweiterungen betreffen Marktvielfalt, Auswertungen zur Gas-Nutzung und
eine klarere Trennung zwischen UI-Vereinfachungen und on-chain-Allgemeingültigkeit.
Die aktuelle Lösung bildet einen didaktisch klaren Kern ab, der technisch
nachvollziehbar bleibt.

\clearpage

% Optionale Literaturliste (benötigt .bib und z.\,B. biblatex/biber).
% \bibliography{references}

\end{document}
